\documentclass[11pt,final,hidelinks]{article}
\usepackage[utf8]{inputenc}
\usepackage[T1]{fontenc}
\usepackage{lmodern}
\usepackage[margin=1in]{geometry}
\usepackage[english]{babel}
\usepackage{graphicx}
\usepackage[activate={true,nocompatibility},final,tracking=true,kerning=true,spacing=true,factor=1100,stretch=10,shrink=10]{microtype}
\microtypecontext{spacing=nonfrench}
\usepackage[noend]{algorithm2e}
\usepackage{caption}
\usepackage{mathtools}
\usepackage{commath}
\usepackage{minted}
\usepackage{enumerate}
\usepackage{amsmath}
\usepackage{subfiles}
\usepackage{booktabs}
\usepackage{amsthm}
\usepackage{amssymb}
\usepackage{pgfplots}
\usepackage{hyperref}
\usepackage{tikz}
\usepackage{tikzscale}
\usepackage[square,numbers]{natbib}
\bibliographystyle{plainnat}
\usepackage{cleveref}
\usepackage[super]{nth}
\usepackage{siunitx}
\usepackage[section]{placeins}
\usepackage{subcaption}

\usepackage{functionnotation}
\usepackage[plain]{setnotation}
\usepackage{approxsetnotation}
\usepackage{approx_data_type_notation}
\usepackage{probabilitynotation}
\usepackage[fancy]{relationnotation}
\usepackage{approxrelationnotation}
\usepackage[section]{envnotation}
\usepackage{algorithmnotation}
%\usepackage{alexmisc}

\newcommand{\amapsto}[2]{\,\ATOverUnder{\mapsto}[#1][\text{\raisebox{3pt}{$#2$}}]\,}

\hypersetup{
    pdftitle={Random approximate values over algebraic types},
    pdfauthor={Alexander Towell},               % author
    pdfsubject={computer science},              % subject of the document
    pdfkeywords={
        probabilistic data structure,
        partial function,
        random approximate set,
        random approximate data type,
        random approximate relation,
        random approximate map,
        random approximate Boolean,
        random approximate unit,
        random approximate sum type,
        random approximate product type,
        random approximate exponential type,
        random approximate void type,
        relation,
        set        
	},                    % keywords
    colorlinks=false,                           % false: boxed links;
    citecolor=green,                            % color of links to 
    filecolor=magenta,                          % color of file links
    urlcolor=green                              % color of external
}

\title
{
    Random approximate values over algebraic types
}
\author
{
    Alexander Towell\\
    \texttt{atowell@siue.edu}
}
\date{}

\begin{document}
\maketitle
\begin{abstract}
We define the semantics of \emph{random approximate values} over algebraic data types: void, unit, Bool, sum (tagged union), product (Cartesian product), and exponential (function space).
From the Bernoulli independence axiom (per-element errors are independent), we derive how approximation distributes over type constructors.
Product types preserve a Kronecker factorization of the joint confusion matrix, so parameters grow additively.
Sum types break this factorization because tag errors couple the components; the optional type $A + \mathrm{Unit}$ recovers the Bernoulli set membership model as a special case.
Exponential types are precisely Bernoulli maps, where the constructive (type-theoretic) and denotational (function-space) perspectives converge.
\end{abstract}

%\microtypesetup{protrusion=false}
\tableofcontents
%\microtypesetup{protrusion=true}


%\listoffigures
%\listofalgorithms

\subfile{sections/intro}
\subfile{sections/approx_values}
%\subfile{sections/bernoulli_type_2}  % subsumed by approx_values.tex
\subfile{sections/prims}
\subfile{sections/composite_types}
\subfile{sections/error_propagation}
\subfile{sections/bernoulli_boolean_algebra}
\subfile{sections/invariants}
\subfile{sections/cpp}

% Orphaned section files (not included above):
%   sections/bernoulli_type_2.tex -- subsumed by sections/approx_values.tex
%   sections/bernoulli_type.tex  -- near-duplicate of cpp.tex (cpp.tex is the more complete version)
%   sections/bernoulli_prolog.tex -- raw plaintext notes on approximate Prolog, not LaTeX

\bibliography{references}
\end{document}
